\documentclass[a4paper]{article}

%% Language and font encodings
\usepackage[english]{babel}
\usepackage[T1]{fontenc}

%% Sets page size and margins
\usepackage[top=3cm,bottom=2cm,left=3cm,right=3cm,marginparwidth=1.75cm]{geometry}

%% Useful packages
\usepackage{amsmath}
\usepackage{amssymb}
\usepackage{graphicx}
\usepackage[colorinlistoftodos]{todonotes}
\usepackage[colorlinks=true, allcolors=blue]{hyperref}


\title{Error Analysis for Solving a Linear System \\ by Cholesky Factorization}
\author{Andrew W. Appel and Pierre Roux\\
with advice from David Bindel}

\begin{document}

\maketitle

\begin{abstract}
  Following Demmel \cite{demmel81} with some lemmas improved by Rump and Jeannerod
\cite{rump14}
  and then implemented in libValidSDP by Roux \cite{DBLP:journals/jar/Roux16}, we derive
  the backward error analysis for solving a s.p.d. linear system $Ax=b$
  with fully accurate consideration of underflows.  We do this as a guide
  to what will be mechanized in Rocq, or what will be documentation of
  what will have been mechanized.
  \end{abstract}

\section{Preliminaries}

\emph{Almost all of this section is quoted directly from Demmel (1981).}

Let $\epsilon$ be the rounding error of the arithmetic, and $\lambda$ be the
underflow threshold.  Suppose, for example, a binary floating point number is
represented as $f\cdot 2^e$ where $f$ lies between 1 and 2 with an $n$ bit fraction,
and $e$ is an integer exponent in the range $e_\mathrm{min} \le e \le e_\mathrm{max}$.
Then $\epsilon=2^{-n}$ (the difference between 1 and the next largest number).
The underflow threshold is $\lambda=2^{e_\mathrm{min}}$; this is the smallest \ldots
normalized number for G.U. \ldots Denormalized numbers lie between 0 and $2^{e_\mathrm{min}}$
in G.U. arithmetic for which the smallest nonzero number is thus $\mu = \epsilon \lambda$.
We will give a more precise model of rounding and underflow errors in the next section.

\newcommand\fl[1]{\ensuremath{\mathit{fl}(#1)}}

Let $\bullet$ be one of the operations $(+,-,\ast,/)$ and let
$\fl{a \bullet b}$ denote the floating point result of the indicated
computation. \ldots To take underflow into account, we write
(\cite{10.1145/1057520.1057522}),
\begin{equation}\tag{11}
\fl{a \bullet b} = (a \bullet b)(1 + e)+\eta\qquad
\mathrm{unless}~a\bullet b ~\mathrm{overflows.}
\end{equation}
In the case of G.U. we have the following constraints on $e$ and $\eta$:
\begin{enumerate}
\item[(1)] $|e|\le \epsilon ~\mathrm{and}~ |\eta|\le \lambda \epsilon$,
\item[(2)] $\eta \times e = 0$
\item[(3)] $\eta=0$ if $\bullet$ is either addition or subtraction.
\end{enumerate}

We define the function $\mathit{UN}(x)$ to be
\[
UN(x)=\left\lbrace\begin{array}{ll}1 & \mathrm{if}~x~\mathrm{underflows}~(x<\lambda)\\
    0 & \mathrm{if~it~does~not}\end{array}\right.
  \]

\newcommand{\norm}[1]{\ensuremath \lVert #1 \rVert}
$\norm{A}$ and $\norm{b}$ denote the norms of matrix $A$ and vector $b$.
$\norm{A}_\infty$ denotes the infinity norm of $A$, and similarly for
$\norm{b}_\infty$.  $k(A)=\norm{A}_\infty\norm{A^{-1}}_\infty$ denotes the condition number
of the matrix $A$.

$|A|~(|b|)$ denotes the matrix (vector) whose entries are the absolute values
  of the corresponding entries in $A$ ($b$).  Inequalities like $|A|>|B|$
  $(|a|>|b|)$ are meant componentwise.

\paragraph{Approach.}
We use backward error analysis.  Thus, when trying
to solve
\begin{equation}\tag{12}
  Ax = b
\end{equation}
for $x$ we obtain an approximation $\hat{x}=x~\delta x$ which satisfies the
perturbed problem
\begin{equation}\tag{13}
  (A+\delta A)\hat{x}=b+\delta b~.
\end{equation}
The task of backwards error analysis is to obtain bounds on $\delta A$ and $\delta b$.
These bounds can be used in turn to bound the residual $r$:
\begin{equation}\tag{14}
  r = A \hat{x}-b = -\delta A \hat{x} + \delta b
\end{equation}

\section{Comparison with Demmel (1981)}

  \paragraph*{Demmel's Proposition 2} is as follows.  As we explain in the
  next sections, we'll use something slightly different.

\newcommand\float[1]{\ensuremath \mathit{float}(#1)}
To analyze solving a triangular system, we need another expression for the error in
computing an inner product:
\begin{equation}\tag{demmel.31}
  \float{\sum_{i=1}^n a_i b_i} = \sum_{i=1}^n a_i b_i (1 + E_i) + \eta ~.
\end{equation}
In the absence of underflow we have
\begin{equation}\tag{demmel.32}
\begin{array}{l}
  |E_1| \le n \epsilon \\
  |E_j| \le (n+2-j)\epsilon \quad \mathrm{if}\quad j>1~.
\end{array}
\end{equation}

In the case of G.U.\footnote{Editor's note: G.U. = Gradual Underflow, that is,
a floating-point format that supports subnormal numbers.  The technical report also covers the case of ``S.Z.'' or ``store zero'', which we omit here in all lemmas and theorems.}
we have the same bounds on the $|E_i|$, and
\begin{equation}\tag{demmel.34}
|\eta| \le n \lambda \epsilon
\end{equation}


\paragraph*{Demmel's Lemma 4: Error Analysis of Solving a Triangular System.}
\emph{As we explain in the next sections, we will use a variation
of this lemma.}

Consider the lower triangular system $Ly=b$ where $L$ does not necessarily
have ones on the diagonal.  (The analyses when only ones are on the
diagonal or L is upper triangular are similar.)  Then in solving
$Ly=b$ using the program in Appendix 2 (where inner products are
accumulated with roundoff error $\epsilon_\mathrm{s}$ and
smallest number $\mu_\mathrm{s}$ we obtain $\hat{y}=y+\delta y$ which
actually satisfies $(L+\delta L)\hat{y}= b + \delta b$, where
\begin{equation}\tag{demmel.80}
  |\delta L_{ij}| \le |L_{ij}| \cdot \left\lbrace
  \begin{array}{lll}
    \epsilon & \mathrm{if} & i=j=1 \\
    (i-1)\epsilon_\mathrm{s} & \mathrm{if} & i>j=1 \\
    (i+1-j)\epsilon_\mathrm{s} & \mathrm{if} & i>j>1 \\
    \epsilon+\epsilon_\mathrm{s} & \mathrm{if} & i=j>1 
  \end{array}\right.
\end{equation}
independent of underflow.

In the absence of underflow
\begin{equation}\tag{demmel.81}
  \delta b = 0~.
\end{equation}
Using G.U.
\begin{equation}\tag{demmel.82}
  \begin{array}{lcl@{~}ll}
    \mathrm{if} & i=1 & |\delta b_i |~ \le &\lambda\epsilon|L_{11}|\mathit{UN}(y_1)
    & \mathrm{if}~y_1~\mathrm{underflows} \\
    \mathrm{if} & i>1 & |\delta b_i |~ \le &i \cdot \lambda_\mathrm{s}\epsilon_\mathrm{s}~ + & \mathrm{intermediate~underflows} \\
    & & & \lambda \epsilon |L_{ii}| \mathit{UN}(y_i) 
    & \mathrm{if}~y_i~\mathrm{underflows}
  \end{array} ~.
\end{equation}
If $L_{ii}$ is known to be 1, the only changes in the above bounds
(besides substituting 1 for $L_{ii}$) are that $E_{11} = \delta b_1 = 0$
independent of underflow.

\section{Error analysis of Cholesky Solve}

Instead of Demmel's Proposition 2, we will be using
the following lemma:

\paragraph*{Rump/Jeannerod Lemma 2.1} (LibValidSDP lemma\_2\_1\_aux): Analysis of $(c-\sum a_i b_i)/b_k$.
  \newline
  Let $\hat{y} = \fl{(c-\sum_{0\le i<k} a_i b_i)/b_k}$ computed by repeated subtraction left to right.  Assume $b_k\not= 0$. Then
\begin{equation}\label{lemma_2_1}\tag{lemma2.1}
  |b_k \hat{y} - (c-\sum_{0 \le i<k}a_i b_i)| \le (k+1)\epsilon \left(|b_k\hat{y}| + \sum_{0\le i < k}|a_ib_i|\right)
  + (1+(k+1)\epsilon)(k+|b_k|)\lambda\epsilon.
  \end{equation}
  
Demmel's Lemma 4 provides fine-tuned bounds for four different cases,
depending on whether $i>j$ and whether $j>0$.
In all these cases we will prove a
weaker bound because that allows us to
re-use the Rump/Jeannerod Lemma 2.1 for $c - \sum x_i y_i$.

\paragraph*{Theorem Demmel\_L4\_adapted:} 
Solving the lower triangular system $Ly = b$ we obtain $\hat{y}=y+\delta y$
which satisfies $(L+\delta L)\hat{y}= b + \delta b$, where
\begin{equation}\tag{demmel\_lemma\_4}
  \begin{array}{rcl}
  |\delta L_{ij}| & \le & n\epsilon|L_{ij}|\qquad\mathrm{and} \\
  |\delta b_i | & \le & (1+n\epsilon)\lambda \epsilon \left(n + |L_{ii}|\right)~.
  \end{array}
\end{equation}
\noindent The factor $(1+n\epsilon)$ was omitted from Demmel's original but is required for correctness.

\paragraph*{Proof.}  From \ref{lemma_2_1}, we get
\begin{equation}
  \left|L_{ii}\hat{y}_i - \left(b_i - \sum_{j=0}^{i-1}L_{ij} \hat{y}_j\right)\right|
  \le (i+1)\epsilon \left(|L_{ii}\hat{y}_i| + \sum_{j=0}^{i-1}|L_{ij}\hat{y}_j|\right)
  + (1 + (i+1)\epsilon)(i + |L_{ii}|)\lambda \epsilon
\end{equation}
that is
\begin{equation}
  \left|b_i - \sum_{j=0}^{i}L_{ij} \hat{y}_j\right|
  \le (i+1)\epsilon \left(\sum_{j=0}^{i}|L_{ij}\hat{y}_j|\right)
  + (1 + (i+1)\epsilon)(i + |L_{ii}|)\lambda \epsilon
\end{equation}
and since $i < n$, we have $i + 1 \le n$ which enables to conclude. Q.E.D.

\section{Solving a linear system by Cholesky+Substitution}

We now derive mixed error bounds for the
floating-point result $\hat{x}$ of solving $Ax=b$ by Cholesky decomposition $A=R^T R$ followed by forward substitution $R^Ty=b$ then back substitution $Rx=y$.

\paragraph*{Lemma th\_2\_3\_aux3: Accuracy of Cholesky decomposition.}
Suppose $A$ is an $n\times n$ symmetric positive definite floating-point matrix,
with $(n+1)\epsilon < 1$.
Let $\hat{R}$ be the result of floating-point Cholesky decomposition (inner-product method)
of  $A$ into $\hat{R}^T \hat{R}$.
Let $\delta A = A - \hat{R}^T \hat{R}$.
Let the vector $\hat{R}_j$ be the $j$th column of $\hat{R}$.
Let $a_\mathrm{max}$ be the maximum diagonal element of $|A|$.
Then
\begin{equation}\tag{2.3aux3}
  |\delta A _{ij}| < (\mathrm{min}(i,j)+2)\epsilon \norm{\hat{R}_i}_2 \norm{\hat{R}_j}_2
  + 4 \lambda\epsilon (n+1+a_\mathrm{max}).
\end{equation}

\paragraph*{Proof.} See Rump+Jeannerod, or see libValidSDP.  Q.E.D.

\paragraph*{Definition.} Define $a'_\mathrm{max}$
as $\frac{a_\mathrm{max}+2(n-1)\lambda \epsilon}{1-(n+1)\epsilon}$.

\paragraph*{Remark.}  Let $A$ and $\hat{R}$ be as above.
Then
\begin{equation}\tag{2.3aux1}
  \norm{\hat{R}_j}_2 \le \sqrt{a'_\mathrm{max}}.
  \end{equation}

\paragraph{Proof.} LibValidSDP   th\_2\_3\_aux1 gives
\begin{equation}
  \norm{\hat{R}_j}_2 \le \sqrt{(A_{jj}+2j\lambda\epsilon)/(1-(j+2)\epsilon)}.
  \end{equation}
and the definition of $a'_\mathrm{max}$ enables to conclude. Q.E.D.

\paragraph{Corollary.} Combining (2.3aux3) and (2.3aux1), we have,
\begin{equation}\tag{ch\_back\_err}
\begin{array}{rl}
|\delta A_{ij}| & <
(\mathrm{min}(i,j)+2)\epsilon
\sqrt{a'_\mathrm{max}} \cdot \sqrt{a'_\mathrm{max}}
+ 4 \lambda\epsilon (n+1+a_\mathrm{max}) \\
& \le
(n+1)\epsilon
a'_\mathrm{max}
+ 4 \lambda\epsilon (n+1+a_\mathrm{max}) \\
\end{array}
\end{equation}

\paragraph*{Lemma.} We obtain bounds on $\delta R^T \hat{R}$ by
combining (2.3aux1) and (demmel\_lemma\_4):
\begin{equation}\tag{csolve\_back\_err\_aux1}
\begin{array}{rl}
|(\delta R^T\hat{R})_{ij}| & \le (|\delta R^T|\cdot|\hat{R}|)_{ij} \\
& \le  \sum_k n\epsilon |\hat{R}_{ki}|\cdot|\hat{R}_{kj}| \qquad \mathrm{(by~demmel\_lemma\_4)}\\
& \le n\epsilon \norm{\hat{R}_i}_2 \norm{\hat{R}_j}_2 \qquad \mathrm{(by~Cauchy-Schwartz)}\\
& \le n\epsilon \sqrt{a'_\mathrm{max}} \sqrt{a'_\mathrm{max}} \quad\qquad\mathrm{(by 2.3aux1)}\\
& = n\epsilon a'_\mathrm{max}.
\end{array}
\end{equation}
The same bound is obtained on $\hat{R}\delta R$.

\paragraph{Lemma.} Bounds on $\delta R^T \delta R$ as follows:
\begin{equation}\tag{csolve\_back\_err\_aux2}
  \begin{array}{rl}
    |(\delta R ^T \delta R)_{ij}| & \le (|\delta R^T| \cdot |\delta R|)_{ij} \\
    & \le \sum_k |\delta R^T|_{ik} |\delta R|_{kj} \\
    & \le \sum_k n^2\epsilon^2 |\hat{R}_{ki}|\cdot|\hat{R}_{kj}| \\
    & \le n^2\epsilon^2 \norm{\hat{R}_i}_2 \norm{\hat{R}_j}_2 \\
    & \le n^2\epsilon^2 a'_\mathrm{max}.
  \end{array}
\end{equation}


\paragraph*{Theorem: Cholesky factor + substitution.}
Suppose $A$ is an $n\times n$ symmetric positive definite floating-point matrix
with minimum eigenvalue $>\sqrt{\lambda}$, with $(n+1)\epsilon < 1$.
Suppose $b$ is a floating-point vector of length $n$.
Suppose $\hat{x}$ is the result of floating-point Cholesky factorization of
$A$ into $\hat{R}^T\hat{R}$ followed by forward substitution finding $\hat{y}$ such
that $\hat{R}^T \hat{y}=b$ and then $\hat{R}\hat{x}=\hat{y}$.
Taking into account underflows in the factorization, but not yet
considering underflows in the substitutions, we have
that there is a backward error $\delta A$ such that
$(A+\delta A)\hat{x}=b$, bounded as follows:
\begin{equation}\tag{csolve\_back\_err}
|\delta A| \le ((3n+1)\epsilon+n^2\epsilon^2)a'_\mathrm{max}+4 \lambda\epsilon (n+1+a_\mathrm{max}).
\end{equation}

\paragraph*{Proof.}
If it were not for the underflow terms $\delta b$ and $\delta y$, we would have
\begin{equation}\tag{xxx}
  \begin{array}{rl}
    b &= (\hat{R}^T+\delta R^T)\hat{y}\\
    & = (\hat{R}^T+\delta R^T)(\hat{R}+\delta R)\hat{x} \\
    & = (\hat{R}^T\hat{R}+\delta R^T\hat{R} +\hat{R}^T\delta R  +\delta R^T \delta R)\hat{x} \\
    & = (A + \delta A' +\delta R^T\hat{R} +\hat{R}^T\delta R  +\delta R^T \delta R)\hat{x} \\
    & = A\hat{x} + (\delta A' +\delta R^T\hat{R} +\hat{R}^T\delta R  +\delta R^T \delta R)\hat{x}. \\
  \end{array}
\end{equation}

We will define 
\begin{equation}\tag{xxx}
  \begin{array}{rl}
    \delta A   & = \delta A' +\delta R^T\hat{R} + \hat{R}^T\delta R  + \delta R^T \delta R\\
      & \le |\delta A'| +|\delta R^T\hat{R}| +|\hat{R}^T\delta R|  +|\delta R^T \delta R|.
  \end{array}
\end{equation}

From (ch\_back\_err) we have
$  |\delta A'_{ij}| < (n+1)\epsilon a'_\mathrm{max}
+ 4 \lambda\epsilon (n+1+a_\mathrm{max}).$\newline
From csolve\_back\_err\_aux1 we have
$|(\delta R^T\hat{R})_{ij}| \le n\epsilon a'_\mathrm{max}$
and similarly 
$|(\hat{R}^T\delta R)_{ij}| \le n\epsilon a'_\mathrm{max}$.\newline
From csolve\_back\_err\_aux2 we have
    $|\delta R^T \delta R|
 \le n^2\epsilon^2 a'_\mathrm{max}$.
\newline
The result follows immediately.  Q.E.D.

\paragraph{Still to do:}  Account for underflows in the forward substitution and back substitution.

\bibliographystyle{alpha}
\bibliography{cholesky}

\end{document}
